% Cote 119: offprints, reports, courses, CV and letters (typed and printed part): archive cote 119, site page 74 (PDF page 75).
% offprint. Draft compared: none.
% AI-transcribed by Claude Opus 5.5 (claude-opus-5-5, adaptive thinking, effort high), Message Batches API, batch r119_full (msgbatch_01DwZNm5Y7YsKBWyAN4sFs8w), 2026-11-13; not yet checked by eye.
\scanpage{119/74}
\mspage{( 2 )}

1.3. Si $C$ est une catégorie ind-tannakienne sur A, A$'$ une A-algèbre,
on appelle foncteur fibre sur $C$ à valeurs dans A$'$ un $\otimes$-foncteur ACU
A$'$-linéaire $\omega : C_{(\mathrm{A}')} \to \mathrm{Mod}(\mathrm{A}')$ qui soit fidèle, exact et commute avec les
limites inductives. En prenant comme morphismes les $\otimes$-morphismes
unifères on obtient une catégorie Fib$(C, \mathrm{A}')$. La collection de ces catégories,
pour A$'$ variable, définit de façon naturelle une catégorie fibrée sur la
catégorie Sch$_{/\mathrm{A}}$ des A-schémas, qu'on notera FIB$(C)$. On remarque alors
que FIB$(C)$ est une gerbe sur Sch$_{/\mathrm{A}}$ pour la topologie fpqc [(${}^{4}$), III, 2.1.1],
et que cette gerbe est liée localement par un groupe affine et plat. Une
gerbe vérifiant cette dernière condition sera appelée tannakienne.

Remarquons enfin que la formation de la gerbe tannakienne FIB$(C)$
est 2-fonctorielle en $C$.

1.4. Soit $\mathcal{G}$ une gerbe tannakienne sur A, i. e. une gerbe sur Sch$_{/\mathrm{A}}$
pour la topologie fpqc, liée localement par un groupe affine et plat. On note
Rep$(\mathcal{G})$ la catégorie des foncteurs cartésiens $\mathcal{G} \to \mathrm{QCOH}(\mathrm{A})$, où QCOH(A)
dénote le champ sur Sch$_{/\mathrm{A}}$ des $\mathcal{O}_{\mathrm{S}'}$-modules quasi-cohérents, pour un A-schéma
variable S$'$. La loi $\otimes$ de QCOH(A) définit sur Rep$(\mathcal{G})$ une loi $\otimes$ ACU pour
laquelle Rep$(\mathcal{G})$ est une catégorie ind-tannakienne sur A. Par exemple,
si G est un A-groupe affine et plat, et $\mathcal{G} = \mathrm{TORS}(\mathrm{G})$ est la gerbe tanna-%
kienne des torseurs à droite sous G, pour la topologie fpqc, on a une équi-%
valence de catégories ind-tannakiennes Rep$(\mathcal{G}) \to \mathrm{Rep}(\mathrm{G})$.

1.5. Soient $C$ une catégorie ind-tannakienne, $\mathcal{G}$ une gerbe tannakienne
sur A. On laisse au lecteur le soin de définir des morphismes canoniques :
\[ C \to \mathrm{Rep}\,(\mathrm{FIB}\,(C), \quad \mathcal{G} \to \mathrm{FIB}\,(\mathrm{Rep}\,(\mathcal{G}))\,. \]

Théorème 1.6. --- Les morphismes précédents sont des équivalences; la
correspondance $C \to \mathrm{FIB}(C)$ définit une 2-anti-équivalence de la 2-catégorie
des catégories ind-tannakiennes sur A avec celle des gerbes tannakiennes sur A,
ayant $\mathcal{G} \to \mathrm{Rep}(\mathcal{G})$ comme quasi-inverse.

2. Le cas d'un corps de base.

2.1. Soit $k$ un corps, $C_0$ une $\otimes$-catégorie ACU abélienne $k$-linéaire
possédant des objets Hom. On dit que $C_0$ est une catégorie tannakienne
sur $k$ s'il existe une extension $k'/k$ et un $\otimes$-foncteur ACU $k$-linéaire
$C_0 \to \mathrm{Mod}(k')$ qui soit fidèle et exact et qui commute avec les Hom. Si $C_0$
est une catégorie tannakienne sur $k$, on définit comme dans [(${}^{9}$), 2.2] la
catégorie fibrée FIB$_0(C_0)$ des foncteurs fibre $C_0 \to \mathrm{Loclib}(\mathrm{T})$, où T est un
$k$-schéma variable.

Proposition 2.2. --- Si $C_0$ est une catégorie tannakienne sur $k$, la $\otimes$-caté-%
gorie $C = \mathrm{Ind}(C_0)$ des ind-objets de $C_0$ [(${}^{5}$), A, 2] est ind-tannakienne
sur $k$; de plus, on a une équivalence de gerbes tannakiennes
\[ (2.2.1) \qquad \mathrm{FIB}\,(C) \xrightarrow{\approx} \mathrm{FIB}_0\,(C_0)\,. \]

\mspage{\hand{74}}
